\documentclass[10pt,a4paper]{amsart}
\usepackage[margin=19mm]{geometry}
\usepackage{amsmath,amssymb,mathtools}
\usepackage[scr=rsfs,cal=euler]{mathalfa}
\usepackage{xfrac}
\usepackage[unicode,hidelinks,bookmarks=false]{hyperref}
\newcommand{\ie}{i.e.\ }
\newcommand{\cf}{cf.\ }
\newcommand{\resp}{respectively\ }
\newcommand{\Subsection}{\subsection}
\providecommand{\nobreakdash}{\nobreak\hbox{-}}
\setlength{\emergencystretch}{2em}
\multlinegap0pt
\usepackage{xcolor}
\usepackage{amssymb}
\usepackage{mathtools}
\newtheorem{proposition}{Proposition}
\newcommand{\R}{\mathbb{R}}
\title{Linearized Attention Cannot Enter the Kernel Regime at Any Practical Width}
\author{Jose Marie Antonio Mi\~noza, Paulo Mario P. Medina, Sebastian C. Iba\~nez}
\date{Source: arXiv:2603.13085v2 (CC BY 4.0)}
\begin{document}
\maketitle

\noindent\emph{Source and licence.} Linearized Attention Cannot Enter the Kernel Regime at Any Practical Width. Version arXiv:2603.13085v2 (CC BY 4.0), distributed by arXiv under the Creative Commons Attribution 4.0 International licence, http://creativecommons.org/licenses/by/4.0/, as recorded in the arXiv licence metadata for this exact version and shown on the abstract page https://arxiv.org/abs/2603.13085v2. Full licence notice: \texttt{licenses/arxiv-2603.13085-CC-BY-4.0.txt}.\par\smallskip

\noindent\textit{Editorial scope.} Counted unit: the article's proposition prop:qkv\_equivalence (source0.tex lines 801--804). The article's complete original proof of it is delivered with it (source0.tex lines 806--830, one \texttt{proof} environment of 25 lines). No mathematics is written, repaired or re-derived: the statement and the proof are the article's own lines, in the article's own notation, and no retained line is substituted or re-flowed.  No further source-local context is needed: every cross-reference of every retained line resolves inside the two delivered spans.  Nothing was omitted from any retained span, so the package carries no omission record.  No retained line contains a citation, so the wrapper carries no bibliography and no bibliography entry.  No retained line contains a figure environment, an image-inclusion command or drawing code, so no image is delivered and the package declares no asset dependency.  Editorial formatting only, in addition: an AMS \texttt{amsart} preamble that restates the article's own declarations and macros used by the retained lines, namely the environments \texttt{proposition} on separate counters, together with the standard packages those lines need.  The retained environments print in this document as Proposition 1 in order of appearance, whereas the article prints the same items with the numbering of its own full document.  The complete native source of the article as distributed in the arXiv e-print for this exact version is bundled unchanged as \texttt{sources/196320/source0.tex}.
\begin{proposition}[Kernel Equivalence Under Linear Projections]
\label{prop:qkv_equivalence}
For full-rank projection matrices $\mathbf{W}_Q, \mathbf{W}_K, \mathbf{W}_V \in \R^{d \times d}$, the generalized linearized attention $f^{\text{att}}_{\mathbf{W}}(\mathbf{X}) = (\mathbf{X}\mathbf{W}_Q)(\mathbf{X}\mathbf{W}_K)^T(\mathbf{X}\mathbf{W}_V)$ induces a kernel with the same structural properties as $K_{\text{LinAttn}}$: data-dependence, transitive similarity propagation, and fourth-order polynomial interactions.
\end{proposition}

\begin{proof}[Proof of Proposition~\ref{prop:qkv_equivalence}]
The generalized linearized attention with projection matrices $\mathbf{W}_Q, \mathbf{W}_K, \mathbf{W}_V \in \R^{d \times d}$ is:
\begin{equation}
f^{\text{att}}_{\mathbf{W}}(\mathbf{X}) = (\mathbf{X}\mathbf{W}_Q)(\mathbf{X}\mathbf{W}_K)^T(\mathbf{X}\mathbf{W}_V)
\end{equation}

For row $i$, this computes:
\begin{align}
[f^{\text{att}}_{\mathbf{W}}(\mathbf{X})]_i &= \sum_{k=1}^n (\mathbf{x}_i^T \mathbf{W}_Q \mathbf{W}_K^T \mathbf{x}_k) (\mathbf{W}_V^T \mathbf{x}_k)
\end{align}
where the attention weight is $(\mathbf{x}_i^T \mathbf{W}_Q)(\mathbf{W}_K^T \mathbf{x}_k)^T = \mathbf{x}_i^T \mathbf{M}_{QK} \mathbf{x}_k$ with $\mathbf{M}_{QK} = \mathbf{W}_Q \mathbf{W}_K^T$.

The induced kernel is the inner product of transformed features:
\begin{align}
K_{\mathbf{W}}(\mathbf{x}_i, \mathbf{x}_j) &= \langle [f^{\text{att}}_{\mathbf{W}}(\mathbf{X})]_i, [f^{\text{att}}_{\mathbf{W}}(\mathbf{X})]_j \rangle \nonumber \\
&= \sum_{k,\ell} (\mathbf{x}_i^T \mathbf{M}_{QK} \mathbf{x}_k)(\mathbf{x}_j^T \mathbf{M}_{QK} \mathbf{x}_\ell)(\mathbf{x}_k^T \mathbf{M}_{VV} \mathbf{x}_\ell)
\end{align}
where $\mathbf{M}_{VV} = \mathbf{W}_V \mathbf{W}_V^T$.

Setting $\mathbf{W}_Q = \mathbf{W}_K = \mathbf{W}_V = \mathbf{I}$ gives $\mathbf{M}_{QK} = \mathbf{M}_{VV} = \mathbf{I}$, recovering exactly:
\begin{equation}
K_{\text{LinAttn}}(\mathbf{x}_i, \mathbf{x}_j) = \sum_{k,\ell} (\mathbf{x}_i^T \mathbf{x}_k)(\mathbf{x}_j^T \mathbf{x}_\ell)(\mathbf{x}_k^T \mathbf{x}_\ell)
\end{equation}
For general full-rank projections, the kernel retains the same structural properties: (i) data-dependence through the (transformed) Gram matrix, (ii) transitive similarity propagation via intermediate summation indices, and (iii) fourth-order polynomial interactions. The projections merely rotate the feature space without altering these structural characteristics.
\end{proof}

\end{document}
